\section{Dinâmica orbital e manobras}
% Para compreender e planejar o movimento das espaçonaves ao redor da Terra e em relação a eles mesmos, é necessário realizar o estudo da dinâmica orbital. 
Nesta seção são apresentados os conceitos que darão suporte às simulações numéricas. Será discorrido sobre as condições de inserção em órbita circular logo após o lançamento, seguindo para as equações do movimento orbital derivadas do problema dos dois corpos.
Em seguida, é apresentado uma explicação sobre a manobra de transferência de Hohmann, utilizada como etapa inicial do trabalho.
Por fim, é apresentada a formulação linearizada do movimento relativo em órbitas circulares, descrita através das equações de Hill/Clohessy--Wiltshire, que permitem modelar a aproximação entre o \emph{chaser} e o \emph{target}.

\subsection{Condições de inserção orbital}
Logo após o lançamento da espaçonave \emph{chaser} ser realizada por um veículo lançador,
% como o \emph{New Glenn} \cite{BlueOrigin_NewGlenn}, \emph{Vulcano} \cite{ULA_VulcanCentaur} ou até o \emph{Falcon 9 Block 5} \cite{SpaceX_Falcon9},
a primeira condição a ser atendida para que ocorra a transferência de órbita, é a inserção em uma órbita estável ao redor da Terra.
% Neste trabalho, a inserção foi considerada ser em uma órbita circular, pois a condição inicial é simplificada e utilizada como referência em missões de RVD/B.

A condição para que a trajetória seja circular é que a força gravitacional forneça a aceleração centrípeta necessária para manter o movimento, portanto, a velocidade de inserção em uma órbita de raio $r$ é dada por
\begin{equation}
    \label{eq:veloc_insercao_orbital}
v_{orb} = \sqrt{\frac{\mu}{r}},
\qquad
\mu = G * M_{Terra}
\end{equation}
onde $\mu$ é o parâmetro gravitacional da Terra, $G$ é a constante de gravitação universal, $M_{Terra}$ é a massa da Terra, incluindo oceanos e atmosfera \cite{IC2021}. 
A equação \eqref{eq:veloc_insercao_orbital} fornece a condição inicial de velocidade que garante a permanência da espaçonave em órbita circular. 
Após estabelecer as condições de inserção orbital, utiliza-se as equações do movimento orbital, derivadas do problema dos dois corpos, detalhadas em \eqref{eq:movimentoTarget}.
\subsection{Equações do movimento orbital}
Assim que o \emph{chaser} estiver inserido em órbita circular, a evolução temporal da espaçonave pode ser descrita pelo problema clássico dos dois corpos, conforme \eqref{sec:movimento_orbital_absoluto}.
O problema dos dois corpos despreza as perturbações externas e considera que a trajetória resultante é exclusivamente dada pela atração gravitacional entre a Terra e o veículo.
Conforme \eqref{eq:eqFundamentalDoisCorpos}, a equação fundamental do movimento orbital é dada por:

\begin{equation}
\ddot {\vec r} = -\frac{\mu}{r^3}\vec r;
\qquad
r = \|\vec r\|,
\end{equation}

em que $\vec{r}$ é o vetor posição da espaçonave em relação ao centro da Terra e $\mu = GM$ é o parâmetro gravitacional do sistema Terra--espaçonave. 
Pode-se caracterizar a órbita através da energia orbital específica, definida como:

\begin{equation}
    \label{eq:energia_orbital}
\varepsilon_{orb} = \frac{v^{2}}{2} - \frac{\mu}{r},
\end{equation}

a qual permanece constante em órbitas keplerianas.
A energia orbital específica é constante para toda a órbita e não dependem de $r$ e nem de $v$ em pontos diferentes.  
Para uma órbita elíptica de semi-eixo maior $a$, tem-se:

\begin{equation}
    \label{eq:energia_orbital_eliptica}
\varepsilon_{elip} = -\frac{\mu}{2a}.
\end{equation}

Ao igualar \eqref{eq:energia_orbital} e \eqref{eq:energia_orbital_eliptica}, obtém-se:

\begin{equation}
\frac{v^{2}}{2} - \frac{\mu}{r} = -\frac{\mu}{2a},
\end{equation}

o que conduz à equação vis-viva:

\begin{equation}
    \label{eq:vis_viva}
v^{2} = \mu \left(\frac{2}{r} - \frac{1}{a}\right).
\end{equation}

Para uma órbita circular de raio $r$, ao substituir a equação
\eqref{eq:veloc_insercao_orbital}, a energia orbital assume o valor constante:

\begin{equation}
\varepsilon_{circ} = -\frac{\mu}{2r}.
\end{equation}

% As equações acima estabelecem a base para a formulação da dinâmica orbital que será empregada nas manobras de transferência de órbita,como a transferência de Hohmann e a modelagem do movimento relativo entre as espaçonaves.

\subsection{Transferência de órbita}
\label{sec:hohmann}
A manobra de dois impulsos, denominada transferência de Hohmann, no contexto de RVD/B, é utilizada para ajustar a altitude orbital do \emph{chaser} em relação à órbita do \emph{target}.
Inicialmente, a espaçonave visitante parte de uma órbita circular de raio $r_1$ até uma órbita circular final de raio $r_2$. Além disso, os impulsos são instantâneos para minimizar o gasto de $\Delta v$ para a transferência.

\begin{figure}[ht]
\centering
\includegraphics[width=0.8\linewidth]{2 Formulacao do problema/Imagens/TransfHohmann.png}
\caption{Transferência de Hohmann. Fonte: Autor, adaptado de MATLAB.}
    \label{fig:transf_hohmann}
\end{figure}

\subsubsection{Descrição da manobra}
Em determinado ponto do perigeu, aplica-se o primeiro impulso, denominado $\Delta v_1$.
A espaçonave irá sair da órbita circular de raio $r_1$ e entrará em uma órbita elíptica de transferência, cujo perigeu está em $r_1$ e o apogeu, em $r_2$. 
% No caso da órbita circular, o perigeu e o apogeu são iguais.
No segundo impulso ($\Delta v_2$), aplicado no apogeu, a órbita é circularizada em $r_2$.
A figura \eqref{fig:transf_hohmann} ilustra o processo.

\subsubsection{Velocidades orbitais}
A velocidade inicial\fn{Apresentadas em $[km/s]$} na órbita circular de raio $r_1$ é similar ao que foi apresentado em \eqref{eq:veloc_insercao_orbital}:
\begin{equation}
    v_1 = \sqrt{\frac{\mu}{r_1}}.
\end{equation}

A velocidade no perigeu da órbita de transferência, de forma similar a \eqref{eq:vis_viva}, é dada por:

\begin{equation}
v_{Transf(1)}
=
\sqrt{2\mu \left(\frac{1}{r_1}
-\frac{1}{2a_{Transf}}\right)},
\end{equation}

onde o semi-eixo maior da órbita de transferência é encontrado conforme:
\begin{equation}
a_{Transf}
=
\frac{r_1 + r_2}{2}
\end{equation}

A variação de velocidade requerida no primeiro impulso é
\begin{equation}
\Delta v_1
=
v_{Transf(1)}
-v_1.
\end{equation}

De forma análoga, a velocidade na órbita circular de raio $r_2$ é
\begin{equation}
v_2
=
\sqrt{\frac{\mu}{r_2}},
\end{equation}

enquanto que a velocidade no apogeu da órbita de transferência é:

\begin{equation}
v_{Transf(2)}
=
\sqrt{2\mu \left(\frac{1}{r_2}
-\frac{1}{2a_{Transf}}\right)}.
\end{equation}

O segundo impulso é então
\begin{equation}
\Delta v_2
=
v_2
-v_{Transf(2)}.
\end{equation}

\subsubsection{Tempo de transferência e custo total}
O tempo total, da transferência orbital corresponde a metade do período da órbita elíptica:

\begin{equation}
t_{Transf}
=
\pi \sqrt{\frac{a_{Transf}^3}{\mu}}.
\end{equation}

Para encontrar o custo de velocidade total da manobra, utiliza-se:

\begin{equation}
    \Delta v_{Hohmann} = \Delta v_1 + \Delta v_2.
\end{equation}

% \subsubsection{Considerações finais}
% Após o término da transferência de Hohmann, o \emph{chaser} encontra-se em uma órbita coplanar e próxima da órbita do \emph{target}, no entanto, estará defasado em fase orbital.
% Se necessário, ajustes podem ser realizados para sincronizar ou reduzir as diferenças entre as anomalias verdadeiras ($\nu$) das espaçonaves.
% Após isso, inicia-se o cálculo da dinâmica do movimento relativo, descrita pelas Equações de Hill/Clohessy--Wiltshire.
\subsection{Equações de Hill/CW}

Como \eqref{eq:relativo_LVLHgeral} estava sendo expresso no quadro do \emph{target}, vetor posição é:
\begin{equation}
\label{eq:posicao_LVLH_target}
\vec r_t
=
\begin{bmatrix}
0 \\
0 \\
-r
\end{bmatrix},
\end{equation}
em que está seguindo a orientação dextrogiro.

E também o vetor de velocidade angular:
\begin{equation}
\label{eq:velocidadeAngular_LVLH_target}
\vec\omega=
\begin{bmatrix}
0 \\
-\omega \\
0
\end{bmatrix}
\end{equation}

Em termos da equação \eqref{eq:relativo_LVLHgeral}, é obtido:

\begin{equation}
\label{eq:velAngular_cross_s}
\vec\omega \times \vec s^* =
\begin{bmatrix}
-\omega z \\
0         \\
\omega x
\end{bmatrix} 
\end{equation}

\begin{equation}
\label{eq:velAngular_cross_cross_s}
\vec\omega \times (\vec\omega \times \vec s^*)
=
\begin{bmatrix}
-\omega^2 x \\
0           \\
-\omega^2 z 
\end{bmatrix}
\end{equation}

\begin{equation}
\label{eq:omega_cross_dsdt}
\vec\omega \times \frac{d^* s^*}{dt}
=
\begin{bmatrix}
-\omega \dot z \\
0 \\
\omega \dot x
\end{bmatrix}
\end{equation}

\begin{equation}
\label{eq:dOmega_cross_s}
\frac{d\vec\omega}{dt} \times \vec s^*
=
\begin{bmatrix}
-\dot\omega z \\
0              \\
\dot\omega x
\end{bmatrix}
\end{equation}

\begin{equation}
\label{eq:jacobiana_LVLH}
\mathbf Js^*
=
\begin{bmatrix}
    1 & 0 & 0 \\
    0 & 1 & 0  \\
    0 & 0 & -2
\end{bmatrix}
\vec s^*
=
\begin{bmatrix}
x \\
y \\
-2z
\end{bmatrix}
\end{equation}

No caso da órbita circular ($e=0$), a velocidade angular
% \footnote{Os autores \cite{fehse2003,sidi1997spacecraft} definem a \emph{velocidade angular} (ou \emph{angular rate} ou até mesmo \emph{angular velocity}) como $\omega$, outros autores, como \cite{bate1971}, a definem como $\eta$} 
$\omega$ é constante e coincide com o movimento médio (\emph{mean motion}, $\eta$), definido por \cite{fehse2003}:

\begin{equation}
\label{eq:omega_constante1}
\vec\omega^2 = \frac{\mu}{r_t^3}
\end{equation}

Ou podendo ser reescrita como:

\begin{equation}
\label{eq:omega_constante2}
\eta^2 = \frac{\mu}{r_t^3}
\end{equation}

Como $r_t=a$, a velocidade angular $\omega^{2}= \eta^2 = \mu/r_t^{3}$, portanto, neste caso em particular de órbita circular, $\omega=n$.
Então, $\dot{\omega}=0$, e ao inseri-lo nos termos da equação \eqref{eq:relativo_LVLHgeral}, é possível obter as equações linearizadas gerais do movimento relativo, conhecidas como Equações de Hill \cite{hill1878,fehse2003, wie2008space}.

\begin{equation}
    \label{eq:equacoes_HillX}
\ddot{x} - 2\omega \dot{z} = \frac{F_x}{m_c}
\end{equation}

\begin{equation}
    \label{eq:equacoes_HillY}
\ddot{y} + \omega^2 y = \frac{F_y}{m_c}
\end{equation}

\begin{equation}
    \label{eq:equacoes_HillZ}
\ddot{z} + 2\omega \dot{x} - 3 \omega^2 z = \frac{F_z}{m_c}
\end{equation}

% Nesta convenção, as equações \eqref{eq:equacoes_HillX}, \eqref{eq:equacoes_HillY} e \eqref{eq:equacoes_HillZ} possuem $x$ é along-track, $y$ cross-track e $z$ radial.

\subsection{Convenção do LVLH}

As equações de Hill apresentadas foram formuladas na convenção do sistema de referência LVLH (\emph{Local Vertical Local Horizontal}), onde os eixos foram definidos da seguinte forma \cite{fehse2003}:
\begin{itemize}
    \item Eixo $x$: tangencial à órbita, na direção do movimento orbital (\emph{along-track}).
    \item Eixo $y$: normal à órbita, completando o sistema ortogonal  (\emph{cross-track}).
    \item Eixo $z$: radial, apontando para o centro da Terra  (\emph{radial}).
\end{itemize}

Neste trabalho, será adotada a convenção de referência utilizada pela NASA, que define os eixos como:

\begin{itemize}
    \item Eixo $x$: radial, apontando do centro da Terra para o satélite (\emph{radial}).
    \item Eixo $y$: tangencial à órbita, na direção do movimento orbital  (\emph{along-track}).
    \item Eixo $z$: normal à órbita, completando o sistema ortogonal  (\emph{cross-track}).
\end{itemize}

Por conta dessa mudança de convenção, é necessário realizar uma reformulação das equações \eqref{eq:equacoes_HillX}, \eqref{eq:equacoes_HillY}, \eqref{eq:equacoes_HillZ}, conforme \cite{sidi1997spacecraft}:

\begin{equation}
    \label{eq:Hill_NASA_X}
    \ddot{x} - 2\omega \dot{y} - 3\omega^2 x = \frac{F_x}{m_c}
\end{equation}

\begin{equation}
    \label{eq:Hill_NASA_Y}
    \ddot{y} + 2\omega \dot{x} = \frac{F_y}{m_c}
\end{equation}

\begin{equation}
    \label{eq:Hill_NASA_Z}
    \ddot{z} + \omega^2 z = \frac{F_z}{m_c}
\end{equation}

As equações em \eqref{eq:Hill_NASA_X} a \eqref{eq:Hill_NASA_Z} evidenciam os termos de Coriolis ($\pm2\omega$ sobre as velocidades cruzadas) e o termo centrípeto ou gravidade--gradiente (multiplicadores de $\omega^2$ sobre $x$ e $z$).
% A convenção adotada, é utilizada pela NASA \cite{Folta2022} em missões espaciais e facilita a integração com modelos dinâmicos e simulações padronizadas.
As equações diferenciais lineares com coeficientes que não variam no tempo, conforme as equações \eqref{eq:equacoes_HillX}, \eqref{eq:equacoes_HillY} e \eqref{eq:equacoes_HillZ}, representam o modelo geral que é válido para uma trajetória relativa arbitrária \cite{debruijn2011,ortolano2017} entre o \emph{chaser} e o \emph{target}, considerando que este último se move sob a influência exclusiva de um campo gravitacional central, ou seja, pode ser considerado uma \emph{espaçonave não-cooperativa} \cite{mahendrakar2023}.

Para reduzir a ordem das matrizes, as equações de Hill foram reescritas na forma de espaço de estados.
% Foi adotada a separação do modelo em dois sistemas independentes: (i) uma para a dinâmica fora do plano, sendo a equação \eqref{eq:equacoes_HillY}, e (ii) outra para a dinâmica no plano, sendo as equações \eqref{eq:equacoes_HillX} e \eqref{eq:equacoes_HillZ}, as equações do movimento.
A forma geral do espaço de estados pode ser expressa como \cite{fehse2003,starek2016,sherrill2014}:

\begin{equation}
    \dot{\mathbf{x}} = \mathbf{A}\,\mathbf{x} + \mathbf{B}\,\mathbf{u}
    \label{eq:state_space_general}
\end{equation}

Onde $\mathbf{A}$ é a matriz de transição, $\mathbf{x}$ é o vetor de estados, $\mathbf{B}$ é a matriz de entradas e $\mathbf{u}$ é o vetor de entradas, que neste caso é a força de controle aplicada no \emph{chaser}.

Para o primeiro caso dentro do plano (\textit{in-plane}), as equações \eqref{eq:equacoes_HillX} a \eqref{eq:equacoes_HillZ}, e considerando o vetor de estados
$[x, y, \dot{x}, \dot{y}]^T$ podem ser escritas conforme:

% \begin{equation}
% \begin{bmatrix}
% \dot{x}(t) \\[0.3em]
% \dot{z}(t) \\[0.3em]
% \ddot{x}(t) \\[0.3em]
% \ddot{z}(t)
% \end{bmatrix}
% =
% \begin{bmatrix}
% 0 & 0 & 1 & 0 \\
% 0 & 0 & 0 & 1 \\
% 0 & 0 & 0 & 2\omega \\
% 0 & 3\omega^2 & -2\omega & 0
% \end{bmatrix}
% \begin{bmatrix}
% x(t) \\[0.3em]
% z(t) \\[0.3em]
% \dot{x}(t) \\[0.3em]
% \dot{z}(t)
% \end{bmatrix}
% +
% \begin{bmatrix}
% 0 & 0 \\
% 0 & 0 \\
% \frac{1}{m_c} & 0 \\
% 0 & \frac{1}{m_c}
% \end{bmatrix}
% \begin{bmatrix}
% F_x \\[0.3em]
% F_z
% \end{bmatrix}
% \label{eq:inplane}
% \end{equation}

% \begin{equation}
%     \label{eq:inplaneXY_NASA}
% \begin{bmatrix}
% \dot{x}(t) \\
% \dot{y}(t) \\
% \ddot{x}(t) \\
% \ddot{y}(t)
% \end{bmatrix}
% =
% \begin{bmatrix}
% 0 & 0 & 1 & 0 \\
% 0 & 0 & 0 & 1 \\
% -3\omega^2 & 0 & 0 & -2\omega \\
% 0 & 0 & 2\omega & 0
% \end{bmatrix}
% \begin{bmatrix}
% x(t) \\
% y(t) \\
% \dot{x}(t) \\
% \dot{y}(t)
% \end{bmatrix}
% +
% \begin{bmatrix}
% 0 & 0 \\
% 0 & 0 \\
% \frac{1}{m_c} & 0 \\
% 0 & \frac{1}{m_c}
% \end{bmatrix}
% \begin{bmatrix}
% F_x \\
% F_y
% \end{bmatrix}
% \end{equation}

\begin{equation}
\label{eq:state_space_xy}
\dot{\vec x}_{xy} =
\underbrace{\begin{bmatrix}
0 & 0 & 1 & 0\\
0 & 0 & 0 & 1\\
3\omega^2 & 0 & 0 & 2\omega\\
0 & 0 & -2\omega & 0
\end{bmatrix}}_{\mathbf A_{xy}}
\vec x_{xy} \;+\;
\underbrace{\begin{bmatrix}
0 & 0\\
0 & 0\\
1/m_c & 0\\
0 & 1/m_c
\end{bmatrix}}_{\mathbf B_{xy}}
\begin{bmatrix} F_x \\ F_y \end{bmatrix},
\qquad
\vec x_{xy}=\begin{bmatrix} x \\ y \\ \dot x \\ \dot y \end{bmatrix}.
\end{equation}

De forma similar, para a dinâmica fora do plano orbital, (\textit{out-of-plane}), considerando o vetor de estados$[z,\dot{z}]^T$, obtém-se:

% \begin{equation}
%     \label{eq:outofplaneZ_NASA}
% \begin{bmatrix}
% \dot{z}(t) \\
% \ddot{z}(t)
% \end{bmatrix}
% =
% \begin{bmatrix}
% 0 & 1 \\
% \omega^2 & 0
% \end{bmatrix}
% \begin{bmatrix}
% z(t) \\
% \dot{z}(t)
% \end{bmatrix}
% +
% \begin{bmatrix}
% 0 \\
% \frac{1}{m_c}
% \end{bmatrix}
% F_z
% \end{equation}

\begin{equation}
\label{eq:state_space_z}
\dot{\vec x}_z =
\underbrace{\begin{bmatrix}
0 & 1\\
-\omega^2 & 0
\end{bmatrix}}_{\mathbf A_z}
\vec x_z \;+\;
\underbrace{\begin{bmatrix}
0\\
1/m_c
\end{bmatrix}}_{\mathbf B_z} F_z,
\qquad
\vec x_z = \begin{bmatrix} z \\ \dot z \end{bmatrix}.
\end{equation}


% \subsubsection{Considerações finais}
% A aproximação relativa será conduzida em R-bar (eixo radial, $x$), sob realimentação de força, enquanto que a orientação do \emph{chaser} é controlada de forma independente para atender restrições de captura.
% Na próxima seção, apresentam-se as equações cinemáticas (quaternions) e dinâmicas (Euler) empregadas na malha de atitude e na definição da orientação relativa $q_{\mathrm{rel}}=q_t^{-1}\otimes q_c$.